\documentclass{jsarticle}
\usepackage[dvipdfmx]{graphicx}
\usepackage{amsmath}
\usepackage{amssymb}
\begin{document}
\title{Explain in Global defferential equation and \\
Global integrate equation. \\
Varintegrate equation, and horizen cut of equations.}
\author{Masaaki Yamaguchi}
\date{}
\maketitle


\newcommand{\variint}{%
	\begin{picture}(15,30)
		\put(0,0){
			$ \iint $}
		\put(0,5){\line(15,0){15}}
	\end{picture} %
}

\newcommand{\variiint}{%
	\begin{picture}(20,30)
		\put(0,0){
			$ \iiint $}
		\put(0,5){\line(15,0){20}}
	\end{picture} %
}



 $$ {\partial \over \partial f}F(x) = \int \int \text{cohom}D_k(x)[I_m]  $$
 Cohomology element is constructed with isotopy of D-brane,
 and this isotopy of symmetry structure is sheaf of manifold,
 more over this brane of element is double integrate of projection.
 $$ {\partial \over \partial f}F = {{}^t \variint} \text{cohom}D_k(x)^{\ll p} $$
 And global partial defferential equation is integrate of cut in cohomolgy,
 and this step of defferential D-brane of sheap value is varint of equals,
 inverse of horizen of cute of section value is reverse of integrate of 
 operators. Three dimension of varint of manifolds in operator of sheaps,
 and this step of times of value is equal with D-brane of equations.
 Global differential equation and integrate equation is operator of 
 global scales of horizen cut and add of cut equation are routs.

 $$ \oint r dx_m = \mathcal{O}(x,y) $$

 $$ {{}^t \variiint}_{{D(\chi,x)}} \text{Hom}[D^2 \psi]^{\ll p} \cong 
 \text{vol}\left({V \over S}\right) $$

 $$ {\partial \over \partial f}F(x) = {{}^t \variint} 
 \text{cohom}D_k(x)^{\ll p} $$

 $$ {d \over df}F = (F)^{f^{'}}, \int F dx_m = (F)^{f}  $$
 $$ {d \over df}\int \int{1 \over (x \log x)^2}dx_m =
 \left(\int \int{1 \over (x \log x)^2}dx_m \right)^{f^{'}}  $$
 $$ f^{'} = (2x (\log x + \log (x + 1))) $$
  $$  \int \int F dx_m = \left({1 \over (x \log x)^2}\right)^{
	  (\int 2x (\log x + \log (x + 1))dx)} $$ 
   $$= e^{-f} $$
  $$ {d \over df}F = \left({1 \over (x \log x)^2}\right)^{
	  (2x (\log x + \log (x + 1))) } $$  
  $$= e^{f} $$
  $$ \log (x \log x) \ge 2(\sqrt{y \log y})^{1 \over 2} $$
 
大域的積分多様体の多重積分は、多様体の階層によって、積分回数が
決まっている。大域的微分多様体は、ニュートン形式とライプニッツ形式の
外微分によって計算される。


  $$ \pi(\chi,x) = [i\pi(\chi,x),f(x)] $$
  $$ \int{1 \over (x \log x)}dx = i \int{1 \over (x \log x)}dx^{N}  
  + {}^{N} i \int{1 \over (x \log x)}dx $$
  $$  \int{1 \over (x \log x)}dx = i \int{x \log x}dx 
  - \int{1 \over (x \log x)}dx $$
  $$ \int \int{1 \over (x \log x)^2}dx_m = i {1 \over 2}x^2 $$
  $$ \int \int{1 \over (x \log x)^2}dx_m = i \int \int_M dx_m $$
  $$ \le {1 \over 2}i + x^2 $$
  $$ E = - {1 \over 2}mv^2 + mc^2 $$
  $$ \lim_{x \to \infty} \int \int{1 \over (x \log x)^2}dx_m \ge {1 \over 2}i $$
  $$ {d \over df}\int \int{1 \over (x \log x)^2}dx_m = {1 \over 2}i $$
  $$ {d \over df}\int \int{1 \over (x \log x)^2}dx_m = 
  \left(\int \int {1 \over (x \log x)^2}dx_m \right)^{(f)^{'}}$$
  $$ (f)^{(f)^{'}} $$
  $$ = (f^f)^{'} $$
  $$ = e^{x \log x} $$

  $$ \Gamma = \int {e^{-x}x^{1-t}}dx, 
  {d \over df}F = e^f, \int F dx_m = e^{-f} $$
  $$ \int x^{1-t}dx = {d \over df}F, \int e^{-x}dx = \int F dx_m $$
  $$ e^{i \theta} = \cos \theta + i \sin \theta $$
  $$ x \bot y \to x \cdot y = \vec{0}, 
  {d \over df}\int \int{1 \over (x \log x)^2}dx_m = {1 \over 2}i $$

  $$ {d \over df}\int \int{1 \over (x \log x)^2}dx_m = 
  \left(\int \int {1 \over (x \log x)^2}dx_m \right)^{(f)^{'}}$$
  $$ {1 \over 2} + {1 \over 2}i = \vec{0}, {1 \over 2}\cdot{1 \over 2}i $$
  $$ A + B  =  \vec{0}, A \cdot B = {1 \over 4}i $$
  $$  \tan 90 \neq 0, ||ds^2|| = A + B $$
  自明な零点は実軸${1 \over 2}  $に1を除いてある。
  $\sin 0 = 0$と$$ e^{i \theta} = \cos \theta + i \sin \theta $$
  $\theta $によって、不確定性原理の関係でもあり、粒子、電子、原子が
  確率分布になっているのも開集合で証明できる。宇宙と異次元の関係にもなっている。
  $$ {d \over df}F = {d \over df} \int \int{1 \over (x \log x)^2}dx_m 
  + {d \over df} \int \int {1 \over (y \log y)^{1 \over 2}}dy_m
  = \Gamma(x,y) $$
  $$ {d \over d \gamma}\Gamma = (e^{-x}x^{1-t})^{\gamma^{'}} \to
  {d \over d \gamma}\Gamma = \Gamma^{(\gamma^{'})}$$

\end{document}
